%%% FIELD cert-statement
For fixed rational \((p_x,r_{0,x},r_{1,x},\Gamma_A,\Gamma_0,\Gamma_1)\) with interior observed cells and finite budgets, \(\mathsf{CERT}\) terminates and returns exact real-algebraic representations, rational isolating intervals, and exact membership certificates for \(\ell_x\) and \(\upsilon_x\).  Its endpoint candidate list is exhaustive: it enumerates every legal analytic outcome-sign branch on the canonical reverse exposure-odds arc and its endpoints, while the forward arc is represented exactly by complementing \(U\); the one-arc reduction rules out every two-dimensional interior optimizer.  It neither admits a cleared-denominator artifact nor omits a legal root, and it handles separately \(\Gamma_A=1\), each \(\Gamma_a=1\), constant-\(U\) corners, root collisions, identically vanishing resultants, and positive-dimensional stationary continua.  This fixed-input output is not a symbolic atlas over varying observed summaries or budgets.
%%% FIELD sharp-statement
For finite budgets and interior observed cells,
\[
 \{\tau(\widetilde P):\widetilde P\in\mathcal M^{\circ}(P_{\mathrm{obs}},\Gamma_A,\Gamma_0,\Gamma_1)\}
 =\{\tau(\widetilde P):\widetilde P\in\mathcal M(P_{\mathrm{obs}},\Gamma_A,\Gamma_0,\Gamma_1)\}
 =\mathcal I_{\mathrm{com}}=[L,R].
\]
The endpoints use one common \((w_{0,x},w_{1,x})\in D_{A,x}\) per stratum.  If \(\Gamma_A>1\) and \((\Gamma_0,\Gamma_1)\neq(1,1)\), every nontrivial stratum endpoint has a canonical-arc optimizer \((w_{0,x},w_{1,x})=(W_{\Gamma_A}(t),t)\) with \(0<t<1\); a forward-arc optimizer is its image under complementing \(U\).  If \(\Gamma_A=1\) or \(\Gamma_0=\Gamma_1=1\), both stratum endpoints equal \(r_{1,x}-r_{0,x}\) and are attained with the strictly positive pair \(w_{0,x}=w_{1,x}=1/2\).  For every selected pair, the endpoint-attaining law uses the feasible conditional Bernoulli potential-outcome marginals and may use any strictly positive coupling of them given \((X,U)\).  The result concerns only the maintained common-binary-confounder class and asserts no equality, inclusion, or numerical-dominance relation with the class in Rosenbaum--Rubin Sections 2--4.
%%% FIELD cert-proof
Fix a stratum \(x\) and suppress its subscript.  Thus \(p,r_0,r_1\in(0,1)\), the three budgets are finite rational numbers at least one, and every conditioning operation below is justified by observed overlap.  The proof uses exact semialgebraic graphing, the proved one-arc reduction, resultant elimination on zero-dimensional strata, signed-remainder root isolation, rational interval Newton validation, and sign-invariant cylindrical algebraic decomposition on exceptional strata.

For an arm \(a\), an endpoint \(\vartheta\in\{\Gamma_a^{-1},\Gamma_a\}\), and a new variable \(z_a\), put
\[
 F_a(w_a,z_a;\vartheta)
 =(1-w_a)(\vartheta-1)z_a^2
 +\{1+(w_a-r_a)(\vartheta-1)\}z_a-r_a
\]
and \(d_a=1+(\vartheta-1)z_a\).  Lemma \(\mathrm{lem{:}binary\text{-}mixture\text{-}inversion}\) proves that
\[
 F_a=0,\qquad 0<z_a<1,\qquad d_a>0                                      \tag{1}
\]
selects exactly one point, namely \(z_a=\kappa(r_a,w_a,\vartheta)\).  Thus adjoining \(z_a\) subject to (1) is an exact semialgebraic graph construction and introduces neither a missing branch nor a spurious quadratic root.  To record simplicity, define
\[
 G_a(z)=(1-w_a)z+w_a\frac{\vartheta z}{1+(\vartheta-1)z}-r_a.
\]
Since \(F_a=d_aG_a\), evaluation on (1) gives
\[
 \frac{\partial F_a}{\partial z_a}
 =d_a\left\{(1-w_a)+\frac{w_a\vartheta}{d_a^2}\right\}>0.              \tag{2}
\]
Every division and implicit differentiation below is therefore legal on the graph.  Moreover,
\[
 m_{a,0}=z_a,\qquad
 m_{a,1}=\frac{\vartheta z_a}{d_a},\qquad
 q_a=r_a+(\pi-w_a)(m_{a,1}-m_{a,0})                                    \tag{3}
\]
is rational with denominators of known positive sign.

We next expose all analytic outcome branches.  Since \(p\in(0,1)\),
\[
 \pi-w_1=(1-p)(w_0-w_1),
 \qquad
 \pi-w_0=p(w_1-w_0).                                                    \tag{4}
\]
Lemma \(\mathrm{lem{:}outcome\text{-}budget\text{-}extrema}\) applied to (3) yields the complete table
\[
\begin{array}{c|c|c}
 &H^-&H^+\\ \hline
 w_1>w_0&(\vartheta_1,\vartheta_0)=(\Gamma_1,\Gamma_0)
         &(\vartheta_1,\vartheta_0)=(\Gamma_1^{-1},\Gamma_0^{-1})\\
 w_1<w_0&(\vartheta_1,\vartheta_0)=(\Gamma_1^{-1},\Gamma_0^{-1})
         &(\vartheta_1,\vartheta_0)=(\Gamma_1,\Gamma_0).
\end{array}                                                             \tag{5}
\]
On \(w_0=w_1\), (4) gives
\[
 H^-=H^+=r_1-r_0.                                                       \tag{6}
\]
If \(\Gamma_a=1\), the two endpoint choices for that arm coincide and the arm contributes no root variable.  Thus (5)--(6), including equality faces, give finitely many and all legal analytic outcome-sign branches.

For \(\Gamma_A>1\), Theorem \(\mathrm{thm{:}one\text{-}arc\text{-}endpoint\text{-}reduction}\) gives
\[
 W_{\Gamma_A}(t)=\frac{\Gamma_A t}{1+(\Gamma_A-1)t},
 \qquad 0\leq t\leq1,                                                  \tag{7}
\]
and proves
\[
 \ell=\min_{0\leq t\leq1}H^-(W_{\Gamma_A}(t),t),
 \qquad
 \upsilon=\max_{0\leq t\leq1}H^+(W_{\Gamma_A}(t),t).                 \tag{8}
\]
The denominator in (7) is at least one.  The same theorem proves strict coordinate monotonicity on both two-dimensional sign cells, so no non-point-identified endpoint has a two-dimensional interior optimizer.  It also proves that complementing \(U\) maps every point of the forward exposure-odds arc to a point of the canonical reverse arc with the identical objective value.  Consequently (7), its two endpoints, and the finite branches (5) are exhaustive for both exposure arcs and for the global endpoint problem; there is no separate two-dimensional stationary system to enumerate.  At \(t=0,1\), the pairs are \((0,0)\) and \((1,1)\), and Proposition \(\mathrm{prop{:}closed\text{-}chart\text{-}extension}\) gives \(q_a=r_a\) without assigning an odds ratio at a null latent cell.  If \(\Gamma_A=1\), the two inequalities defining \(D_A\) imply \(w_0=w_1\); (6) then gives \(\ell=\upsilon=r_1-r_0\).  The same equality follows when \(\Gamma_0=\Gamma_1=1\).  These point-identified cases are evaluated directly.

Consider one nonconstant analytic branch after substituting (7), and write its rational objective as \(\Phi(t,z_0,z_1)\), constrained by \(F_0=F_1=0\); omit \(z_a\) whenever \(\Gamma_a=1\).  By (2), implicit differentiation is legal, and a tangential stationary point satisfies
\[
 F_0=F_1=0,
 \qquad
 \frac{\partial\Phi}{\partial t}
 -\sum_{a=0}^1
   \frac{\partial\Phi}{\partial z_a}
   \frac{\partial F_a/\partial t}{\partial F_a/\partial z_a}=0.        \tag{9}
\]
Multiplying by the positive factors in (1)--(3), (7), and (2) turns (9) into a polynomial system without changing its legal solutions.  The branch endpoints and every equality face have already been assigned separate cells.  Saturating the stationary ideal by the cleared mixture denominators, the arc denominator, the legal-root derivatives, and the strict branch-boundary polynomials removes only components on which the rational derivation is undefined or which belong to those separate lower-dimensional cells.  Because every saturated factor is nonzero on the current cell, saturation removes no legal point from that cell.

Suppose first that a saturated stationary system is zero-dimensional.  Successive resultants eliminate the legal root variables.  At each elimination \(\mathsf{CERT}\) separately records the vanishing leading-coefficient and subresultant cases.  Hence the elementary implication
\[
 \text{common finite zero}\quad\Longrightarrow\quad\text{resultant zero} \tag{10}
\]
shows that no solution is lost.  Euclidean polynomial division terminates because degree strictly decreases.  Square-free factorization leaves finitely many real roots.  For every square-free univariate factor, the deterministic signed-remainder chain has constant sign variation away from its roots, and the variation drops by one at each simple root; therefore the difference of endpoint variations counts the roots in any rational interval.  Repeated rational bisection produces disjoint rational isolating intervals for every projected root.  Subresultant back-substitution and the inequalities (1) lift those intervals to all and only the real solutions of the original stationary system.

For a nonsingular lifted solution, rational interval Newton contraction gives a box mapped strictly into itself and thereby certifies a unique zero.  If projected roots collide or a lifted isolated solution is singular, gcd, square-free, subresultant, and Thom-sign data record the collision exactly; interval Newton flags rather than deletes it, and exact algebraic sign evaluation against the original equations and inequalities supplies its membership certificate.  In every case the original uncleared equations, positive-denominator conditions, legal-root inequalities, branch conditions, and \(0\leq t\leq1\) are checked.  Thus a wrong quadratic root or a cleared-denominator artifact is rejected.

If a resultant vanishes identically or the stationary ideal is positive-dimensional, \(\mathsf{CERT}\) invokes sign-invariant cylindrical algebraic decomposition of the finite polynomial family consisting of the original constraints, branch boundaries, derivative equations, and objective comparisons.  Termination follows by induction on the number of variables.  In one variable the signed-remainder isolation above produces finitely many points and intervals.  At each projection step, only finitely many coefficients, discriminants, and pairwise resultants are added.  On every lower-dimensional cell where those projection polynomials have fixed nonzero signs, degrees and multiplicities are fixed and distinct real roots cannot cross, so finitely many ordered sections and sectors lift the decomposition.  On a cell where a projection polynomial vanishes, specialization lowers a degree or records an equality and the same finite recursion applies.  The variable count falls at projection and every polynomial family remains finite.  Hence the recursion terminates and returns a stationary continuum or legal-root collision as an exact algebraic cell instead of discarding it.

Exhaustiveness is now immediate but we state it carefully.  Proposition \(\mathrm{prop{:}closed\text{-}chart\text{-}extension}\) makes each objective continuous on the compact exposure region, so its extrema exist.  Theorem \(\mathrm{thm{:}one\text{-}arc\text{-}endpoint\text{-}reduction}\) places every non-point-identified optimizer on the canonical arc, up to the complement involution.  The finite stratification by (5)--(7) places it either at a listed zero-dimensional endpoint or in the relative interior of a one-dimensional analytic cell.  In the latter case its tangential derivative vanishes unless the objective is constant on the cell; isolated zeros are retained by the resultant branch and a nonisolated zero set is retained by cylindrical decomposition.  The point-identified strata were evaluated directly.  Thus every optimizer is represented by an enumerated algebraic point or cell.

Finally, algebraic objective values at distinct candidates are compared exactly by refining their rational isolating intervals until disjoint.  If two values coincide, reduction in their common algebraic coordinate ring certifies equality.  Exact sign evaluation supplies membership certificates for all original constraints.  Performing these finite operations for \(H^-\) and \(H^+\) returns exact real-algebraic representations, rational isolating intervals, and exact membership certificates for \(\ell_x\) and \(\upsilon_x\).  Repeating over the finite set \(\mathcal X\) proves the theorem.  All coefficients were specialized fixed rational inputs; no step returns a symbolic atlas over varying inputs.
%%% FIELD strict-proof
For a stratum \(x\), abbreviate a posterior pair by \(d\in D_{A,x}\).  For every such \(d\), elementary minimization and maximization give
\[
 q^-_{1,x}(d)-q^+_{0,x}(d)
 \geq \min_{e\in D_{A,x}}q^-_{1,x}(e)
      -\max_{f\in D_{A,x}}q^+_{0,x}(f).                                \tag{11}
\]
Taking the minimum of the left side, multiplying by \(\lambda_x>0\), and summing proves \(L\geq L_{\mathrm{sep}}\).  Similarly,
\[
 q^+_{1,x}(d)-q^-_{0,x}(d)
 \leq \max_{e\in D_{A,x}}q^+_{1,x}(e)
      -\min_{f\in D_{A,x}}q^-_{0,x}(f),                                \tag{12}
\]
so \(R\leq R_{\mathrm{sep}}\).  Theorem \(\mathrm{thm{:}sharp\text{-}common\text{-}u\text{-}envelope}\) identifies the common-pair interval and Definition \(\mathrm{def{:}separate\text{-}pair\text{-}relaxation}\) defines the separate-pair outer relaxation.  Hence
\[
 L_{\mathrm{sep}}\leq L\leq R\leq R_{\mathrm{sep}},
 \qquad \mathcal I_{\mathrm{com}}\subseteq\mathcal I_{\mathrm{sep}}.  \tag{13}
\]

We verify the law in Definition \(\mathrm{def{:}rational\text{-}witness}\).  Its marginal confounder prevalence and treatment odds ratio are
\[
 \pi=\frac12\frac45+\frac12\frac15=\frac12,
 \qquad
 \operatorname{OR}\!\left(\frac45,\frac15\right)=16.                 \tag{14}
\]
Its observed risks and outcome odds ratios are
\[
\begin{aligned}
 r_1&=\frac45\frac{13}{20}+\frac15\frac1{20}=\frac{53}{100},
 &\theta_1&=\frac{(13/20)(19/20)}{(1/20)(7/20)}=\frac{247}{7}<36,\\
 r_0&=\frac15\frac45+\frac45\frac1{10}=\frac6{25},
 &\theta_0&=\frac{(4/5)(9/10)}{(1/10)(1/5)}=36.
\end{aligned}                                                           \tag{15}
\]
Bayes' rule gives \(P^\star(A=1\mid U=1)=4/5\) and \(P^\star(A=1\mid U=0)=1/5\).  These probabilities and all four outcome risks are strictly between zero and one.  The conditional independent-Bernoulli construction in Definition \(\mathrm{def{:}rational\text{-}witness}\) therefore makes every full-data cell positive and gives \((Y(0),Y(1))\perp A\mid U\).  Equations (14)--(15) verify the four sensitivity inequalities.  Thus \(P^\star\in\mathcal M^{\circ}\), and
\[
 \tau(P^\star)
 =\frac12\left(\frac{13}{20}+\frac1{20}-\frac45-\frac1{10}\right)
 =-\frac1{10},
 \qquad r_1-r_0=\frac{29}{100}.                                        \tag{16}
\]

It remains to validate the endpoint transcript rather than treating its construction as proof.  On the canonical reverse arc in (C8), \(w_0=W(t)>t=w_1\), and
\[
 w_0(1-w_1)=16w_1(1-w_0),
 \qquad
 16w_0(1-w_1)-w_1(1-w_0)=255w_1(1-w_0)>0.                              \tag{17}
\]
Since \(p=1/2\), \(\pi-w_1=(W-t)/2>0\) and \(\pi-w_0=(t-W)/2<0\).  Lemma \(\mathrm{lem{:}outcome\text{-}budget\text{-}extrema}\) makes the common lower branch \(J_{1/36}\) and upper branch \(J_{36}\), exactly as recorded in (C1)--(C2).  Theorem \(\mathrm{thm{:}one\text{-}arc\text{-}endpoint\text{-}reduction}\) proves that neither common endpoint has a two-dimensional interior optimizer and that complementing \(U\) duplicates the forward-arc values.  Hence (C8) reduces the common optimization exhaustively to this one arc, its diagonal, and its constant-\(U\) endpoints.  The same theorem's strict derivative inequalities applied to a single adjustment term show that each separate armwise extremum also enlarges strictly with posterior separation, so its optimizer is on an exposure boundary; complementing \(U\) again reduces it to the canonical arc.

For the one-dimensional common candidates, (C2) is the numerator obtained by differentiating the two mixture equations along the arc.  Equation (2) in the proof of Theorem \(\mathrm{thm{:}complete\text{-}semialgebraic\text{-}certificate}\) shows that neither legal-root derivative vanishes.  Thus every legal stationary point must annul \(g_-(t)\) or \(g_+(t)\) in (C3)--(C4).  The deterministic Sturm recurrence specified before (C5), together with the variation differences \(8-6=2\), proves that each degree-sixteen polynomial has exactly two roots in \((0,1)\).  The four one-unit variation drops in (C5) isolate all four roots.  Substitution into the original three-equation systems, the nonzero Jacobian enclosures (C7), the legal-risk and denominator checks following (C7), and rational interval Newton contraction validate exactly one legal stationary triple in each box in (C6).  Thus no resultant root is missed, duplicated, or admitted through denominator clearing.

The diagonal and both corner objectives equal \(29/100\).  Comparing this value and all four objective enclosures in (C6) shows that the first lower row and the second upper row are the selected extrema and gives
\[
 L\in[-107500302/10^9,-107500301/10^9],
 \qquad
 R\in[539509967/10^9,539509968/10^9].                                  \tag{18}
\]
The complete-certificate theorem justifies the exact saturation, resultant, Sturm, interval-Newton, and membership operations, while the displayed (C1)--(C8) transcript supplies the complete fixed rational input, cell inventory, stationary polynomials, isolations, and competitor comparisons.

For the separate extrema, (C9) is obtained by the same exact implicit differentiation and elimination.  The extra resultant root \(t=9/325\) in \((0,1)\) forces \(z_0=414/395>1\) and is rejected by Lemma \(\mathrm{lem{:}binary\text{-}mixture\text{-}inversion}\); the other recorded extra roots are either outside \((0,1)\) or have the displayed infeasible root values.  Direct rational substitution of the four legal records in (C10), followed by comparison with the common endpoint value at the arc endpoints, yields
\[
\begin{aligned}
 L_{\mathrm{sep}}
 &=\frac{94711}{277400}-\frac{2553}{5650}
 =-\frac{3461701}{31346200}
 \in[-110434471/10^9,-110434470/10^9],\\
 R_{\mathrm{sep}}
 &=\frac{212689}{302600}-\frac{1797}{12475}
 =\frac{84380923}{150997400}
 \in[558823681/10^9,558823682/10^9].
\end{aligned}                                                           \tag{19}
\]
Exact interval subtraction of (18)--(19) gives
\[
 L-L_{\mathrm{sep}}
 \in[2934168/10^9,2934170/10^9],
 \qquad
 R_{\mathrm{sep}}-R
 \in[19313713/10^9,19313715/10^9].                                     \tag{20}
\]
Both endpoint inclusions in (13) are therefore strict at the witness.

Finally, we prove relative-open stability without assuming persistence of an active branch.  Let \(\gamma\geq1\) vary near \(16\).  On the interior square, the coordinatewise logit map writes the exposure region as \(|s_1-s_0|\leq\log\gamma\).  If \(\gamma_n\to\gamma\), keep \((s_0+s_1)/2\) fixed and truncate \(s_1-s_0\) to \([-log\gamma_n,\log\gamma_n]\).  The inverse-logit pair belongs to \(D_{\gamma_n}\) and converges to the original pair.  The two constant corners belong to every \(D_\gamma\); conversely, every limit of pairs in \(D_{\gamma_n}\) satisfies the two closed cross-product inequalities for \(D_\gamma\).  Therefore
\[
 d_H(D_{\gamma_n},D_\gamma)\longrightarrow0.                           \tag{21}
\]

Lemma \(\mathrm{lem{:}binary\text{-}mixture\text{-}inversion}\) and Proposition \(\mathrm{prop{:}closed\text{-}chart\text{-}extension}\) make all endpoint objectives jointly continuous in the interior observed summaries, finite positive budgets, and posterior pair.  Given minimizers in \(D_{\gamma_n}\), compactness supplies subsequential limits, (21) makes them feasible in \(D_\gamma\), and joint continuity gives the lower-limit inequality for the minima.  Approximating a minimizer in \(D_\gamma\) by points in \(D_{\gamma_n}\) gives the upper-limit inequality.  Maxima follow by negation, and the same argument applies to all four armwise separate extrema.  Thus \(L-L_{\mathrm{sep}}\) and \(R_{\mathrm{sep}}-R\) are continuous at the witness.  Their positive lower bounds in (20) remain positive on a sufficiently small relatively open neighborhood of the displayed interior law and finite budgets.  This proves stable strict tightening.
%%% FIELD atlas-proof
The proof uses first-order semialgebraic elimination to construct the atlas and compact-set arguments for monotonicity and continuity.  All observed summaries are real with \(0<p_x,r_{a,x}<1\), and every budget lies in the compact finite box \(\mathcal B\subset[1,\infty)^3\).

For each stratum, Theorem \(\mathrm{thm{:}one\text{-}arc\text{-}endpoint\text{-}reduction}\) represents every non-point-identified endpoint on
\[
 (w_0,w_1)=\left(\frac{\Gamma_A t}{1+(\Gamma_A-1)t},t\right),
 \qquad 0\leq t\leq1,                                                   \tag{22}
\]
and represents the forward arc by complementing \(U\).  Theorem \(\mathrm{thm{:}complete\text{-}semialgebraic\text{-}certificate}\) shows that every legal root, analytic outcome branch, stationary condition, endpoint condition, and objective comparison in (22) has a finite polynomial description after adjoining legal mixture-root variables.  Reciprocal endpoints are encoded by \(\Gamma_a\vartheta_a=1\), legal because \(\Gamma_a\geq1\).  The strata \(\Gamma_A=1\), \(\Gamma_a=1\), \(t=0,1\), root collisions, and vanishing denominators are recorded separately.

Let \(\mathsf{Feas}_x(v,s;\gamma)\) be the resulting finite disjunction asserting that \(v\) contains legal canonical-arc and mixture-root variables and that the corresponding lower objective equals \(s\), where \(\gamma=(\Gamma_A,\Gamma_0,\Gamma_1)\).  The lower endpoint graph is exactly
\[
\begin{split}
 \operatorname{Graph}(\ell_x)=\{(\gamma,s):{}&
 \exists v\,\mathsf{Feas}_x(v,s;\gamma)\ \wedge\\
 &\neg\exists(v',s')\,[\mathsf{Feas}_x(v',s';\gamma)\wedge s'<s]\}.
\end{split}                                                             \tag{23}
\]
Compactness of \([0,1]\) and continuity of the legal-root objectives ensure that (23) contains an attained minimum, so it neither loses an infimum nor introduces a false endpoint.

Apply the finite projection-and-lifting construction proved in the certificate theorem to (23), treating the observed summaries and \(\gamma\) as coefficients.  Projection eliminates all quantified coordinates and partitions \(\mathcal B\) into finitely many sign-invariant semialgebraic cells over \(K_{\mathrm{in}}\).  On each cell the legal optimizer strata and their algebraic objective values are fixed.  Adjoin variables \((s_x)_{x\in\mathcal X}\), the graph equations (23), and
\[
 L=\sum_{x\in\mathcal X}\lambda_xs_x,
 \qquad L=0.                                                            \tag{24}
\]
Projection and lifting applied to this finite polynomial family give a finite semialgebraic stratification of \(\mathfrak T_0(P_{\mathrm{obs}})\cap\mathcal B\).  Adding branch-boundary polynomials, discriminants, legal-root Jacobian minors, and pairwise objective differences to the projection family records optimizer ties, legal-root collisions, constant-\(U\) contacts, and stationary continua as designated lower-dimensional cells in the joint parameter--optimizer stratification.  Such sets are retained rather than discarded by a generic resultant.

We next prove coordinatewise monotonicity.  By Lemma \(\mathrm{lem{:}outcome\text{-}budget\text{-}extrema}\),
\[
 q^-_{a,x}(w;\Gamma_a)
 =\min_{\theta\in[\Gamma_a^{-1},\Gamma_a]}q_{a,x}(w;\theta),
 \qquad
 q^+_{a,x}(w;\Gamma_a)
 =\max_{\theta\in[\Gamma_a^{-1},\Gamma_a]}q_{a,x}(w;\theta).           \tag{25}
\]
If \(\Gamma_1\) increases, the first interval expands, so \(q^-_{1,x}\) cannot increase.  If \(\Gamma_0\) increases, the second interval expands, so \(q^+_{0,x}\) cannot decrease.  Therefore \(H_x^-=q^-_{1,x}-q^+_{0,x}\) cannot increase under either enlargement.  If \(\Gamma_A'\geq\Gamma_A\), each cross-product inequality defining \(D_{A,x}(\Gamma_A)\) remains true at \(\Gamma_A'\), whence
\[
 D_{A,x}(\Gamma_A)\subseteq D_{A,x}(\Gamma_A').                        \tag{26}
\]
Minimizing the pointwise nonincreasing objective over the enlarged set proves that every \(\ell_x\), and hence \(L=\sum_x\lambda_x\ell_x\) with \(\lambda_x>0\), is coordinatewise nonincreasing.

For continuity, Lemma \(\mathrm{lem{:}binary\text{-}mixture\text{-}inversion}\) and Proposition \(\mathrm{prop{:}closed\text{-}chart\text{-}extension}\) make \(q_{a,x}\) jointly continuous in the posterior pair and every finite positive tilt.  Taking the minimum or maximum over the compact interval in (25) makes \(q^-_{a,x}\) and \(q^+_{a,x}\) jointly continuous in the posterior pair and finite outcome budget.

It remains to control the moving exposure region.  Put \(c=\log\Gamma_A\).  In the interior square, with \(s_a=\operatorname{logit}(w_a)\), the two exposure inequalities are equivalent to
\[
 |s_1-s_0|\leq c.                                                       \tag{27}
\]
If \(c_n\to c\), keep \((s_0+s_1)/2\) fixed and truncate \(s_1-s_0\) to \([ -c_n,c_n]\); the inverse-logit pair lies in \(D_{A,x}(e^{c_n})\) and converges to the original pair.  The constant corners belong to every region.  Conversely, every convergent sequence of feasible pairs has a feasible limit because the cross-product inequalities are closed.  Thus \(D_{A,x}(e^{c_n})\) converges in Hausdorff distance to \(D_{A,x}(e^c)\), including \(c=0\), where the limit is the diagonal.

On the compact product of the posterior square and a finite budget neighborhood, the lower objective is uniformly continuous.  Evaluating the two minima at nearest feasible points in the two Hausdorff-close regions bounds their difference by that uniform-continuity modulus.  Therefore each \(\ell_x\), and hence the finite weighted sum \(L\), is continuous.

If the summaries and box endpoints are rational or real algebraic, every projection coefficient is an effective element of \(K_{\mathrm{in}}\).  Euclidean remainder sequences, exact algebraic-number sign tests, and the finite projection-and-lifting recursion terminate and return exact quantifier-free sign cells and algebraic optimizer strata.  Signed-remainder isolation returns rational isolating intervals after a zero-dimensional specialization.  A positive-dimensional stratum is returned by exact sign conditions rather than by a purported isolating interval.  This is exactly the contract of \(\mathsf{ATLAS}\), distinct from the fixed-input contract of \(\mathsf{CERT}\).
%%% FIELD inference-proof
The proof uses a local analytic implicit expansion, the stated asymptotically linear representation, and a one-dimensional root expansion.  Strictly positive observed cells put \(b_0\) in the relative interior of its \((4J-1)\)-dimensional chart.  The assumption that \(\gamma_0\) is interior to \(\mathcal B\cap(1,\infty)^3\) permits two-sided perturbations of every budget coordinate.  Nonzero denominators and strict inactive inequalities persist by continuity.  The unique active branch and strict objective separation therefore give, on a sufficiently small neighborhood,
\[
 F\{z(b,\gamma),b,\gamma\}=0,
 \qquad
 L(b,\gamma)=\phi\{z(b,\gamma),b,\gamma\}.                             \tag{28}
\]

Let \(J=F_z(z_0,b_0,\gamma_0)\), which is invertible by hypothesis, write \(h=(h_b,h_\gamma)\), and put \(v=z-z_0\).  Differentiability gives
\[
 F(z_0+v,b_0+h_b,\gamma_0+h_\gamma)
 =Jv+F_bh_b+F_\gamma h_\gamma+\rho(v,h),                               \tag{29}
\]
where \(\rho(0,h)=o(\|h\|)\) and
\[
 \|\rho(v,h)-\rho(v',h)\|=o(1)\|v-v'\|
\]
uniformly as \((v,v',h)\to0\).  On a sufficiently small fixed neighborhood, the map
\[
 v\longmapsto-J^{-1}\{F_bh_b+F_\gamma h_\gamma+\rho(v,h)\}
\]
has Lipschitz constant below one and maps a ball of radius \(C\|h\|\) into itself for a sufficiently large constant \(C\).  Its iterates are Cauchy, because successive distances decrease geometrically, and their limit is the unique solution of (29) in that ball.  Hence \(v=O(\|h\|)\), and substitution in (29) gives
\[
 z(b_0+h_b,\gamma_0+h_\gamma)-z_0
 =-J^{-1}F_bh_b-J^{-1}F_\gamma h_\gamma+o(\|h\|).                      \tag{30}
\]

Expanding the second identity in (28) and substituting (30) gives
\[
\begin{aligned}
 L(b_0+h_b,\gamma_0+h_\gamma)-L(b_0,\gamma_0)
 &=\phi_z\{z(b_0+h_b,\gamma_0+h_\gamma)-z_0\}
   +\phi_bh_b+\phi_\gamma h_\gamma+o(\|h\|)\\
 &=\bigl(\phi_b-\phi_zF_z^{-1}F_b\bigr)h_b
   +\bigl(\phi_\gamma-\phi_zF_z^{-1}F_\gamma\bigr)h_\gamma
   +o(\|h\|)\\
 &=\dot L_bh_b+\dot L_\gamma h_\gamma+o(\|h\|).
\end{aligned}                                                           \tag{31}
\]
This proves the displayed derivative formulas.

The assumed asymptotically linear representation implies \(\widehat b_n-b_0=O_P(n^{-1/2})\).  Take \(h_b=\widehat b_n-b_0\), \(h_\gamma=0\), and use \(L(b_0,\gamma_0)=0\) in (31).  Then
\[
 \sqrt n\,L(\widehat b_n,\gamma_0)
 =\dot L_b\sqrt n(\widehat b_n-b_0)+o_P(1)
 \ \Longrightarrow\ \dot L_bZ.                                      \tag{32}
\]
Since \(Z\sim N(0,\Sigma)\),
\[
 \dot L_bZ\sim N(0,\sigma_L^2),
 \qquad \sigma_L^2=\dot L_b\Sigma\dot L_b^{\mathsf T}.               \tag{33}
\]

For the frontier displacement, define
\[
 \nu_0=\frac{\dot L_\gamma^{\mathsf T}}{\|\dot L_\gamma\|}.
\]
The nonzero-gradient assumption gives
\[
 \dot L_\gamma\nu_0=\|\dot L_\gamma\|=:c>0.                         \tag{34}
\]
By continuity of the branch derivative, \(t\mapsto L(b,\gamma_0+t\nu_0)\) has derivative at least \(c/2\) for all sufficiently small \((b-b_0,t)\).  Choose deterministic \(M_n\to\infty\) with \(M_n/\sqrt n\to0\).  Equation (32) implies that, with probability tending to one,
\[
 |L(\widehat b_n,\gamma_0)|<\frac{cM_n}{4\sqrt n}.
\]
The mean-value inequality makes the values at \(t=-M_n/\sqrt n\) and \(t=M_n/\sqrt n\) have opposite signs.  Continuity and strict monotonicity therefore give the unique local root \(t_n\), and a second application of the mean-value inequality gives
\[
 |t_n|\leq\frac{2}{c}|L(\widehat b_n,\gamma_0)|=O_P(n^{-1/2}).          \tag{35}
\]

Use \(h_b=\widehat b_n-b_0\) and \(h_\gamma=t_n\nu_0\) in (31).  Equations (34)--(35) make the remainder \(o_P(n^{-1/2})\), so the root equation becomes
\[
 0=\dot L_b(\widehat b_n-b_0)+\|\dot L_\gamma\|t_n+o_P(n^{-1/2}).     \tag{36}
\]
Multiplication by \(\sqrt n\), solution for \(\sqrt n t_n\), and (32) yield
\[
 \sqrt n\,t_n
 \ \Longrightarrow\
 -\frac{\dot L_bZ}{\|\dot L_\gamma\|}
 \sim N\!\left(0,\frac{\sigma_L^2}{\|\dot L_\gamma\|^2}\right).    \tag{37}
\]
Equation (31) identifies \(\dot L_\gamma=\nabla_\gamma L(b_0,\gamma_0)\), so the standard-error constant is \(\sigma_L/\|\nabla_\gamma L\|\).  If \(\sigma_L>0\), consistency of the stated plug-in estimators and continuity of matrix multiplication, the Euclidean norm, and division by the nonzero denominator give the usual pointwise normal interval.  Replacing the active minimizing branch and \(L\) by the active maximizing branch and \(R\) repeats (28)--(37), proving the final assertion.
%%% FIELD sharp-proof
Define the observed-law identified set
\[
 \Theta_I(P_{\mathrm{obs}})
 =\{\tau(\widetilde P):
       \widetilde P\in\mathcal M(P_{\mathrm{obs}},\Gamma_A,\Gamma_0,\Gamma_1)\}.
                                                                            \tag{38}
\]
Here \(P_{\mathrm{obs}}=\mathcal L(X,A,Y)\) is a law on \(\mathcal X\times\{0,1\}^2\), \(\widetilde P\) is a law on \(\mathcal X\times\{0,1\}^4\), and \(\tau(\widetilde P)=\widetilde E\{Y(1)-Y(0)\}\).  We use the declared finite support, \(\lambda_x>0\), observed overlap \(0<p_x<1\), interior risks \(0<r_{a,x}<1\), and finite budgets \(1\leq\Gamma_A,\Gamma_0,\Gamma_1<\infty\).  These conditions justify conditioning and division and keep every finite tilt denominator positive.  No rank or matrix-inversion condition is used.

Fix \(x\in\mathcal X\) and first take \(\widetilde P\in\mathcal M^{\circ}\).  Assumption \(\mathrm{ass{:}observed\text{-}margin}\) identifies \(\lambda_x,p_x,r_{a,x}\) with the full-law summaries.  Assumption \(\mathrm{ass{:}consistency}\) gives
\[
 r_{a,x}=\widetilde P\{Y(a)=1\mid A=a,X=x\}.
\]
Assumption \(\mathrm{ass{:}latent\text{-}ignorability}\), followed by conditioning on binary \(U\), yields
\[
 r_{a,x}=(1-w_{a,x})\mu_{a,0,x}+w_{a,x}\mu_{a,1,x}.                    \tag{39}
\]
Bayes' rule is legal because every full-data cell is positive and gives
\[
 \pi_x=\widetilde P(U=1\mid X=x)
      =p_xw_{1,x}+(1-p_x)w_{0,x}.                                      \tag{40}
\]
The four positive \((A,U,X=x)\)-cell probabilities give
\[
 \frac{\zeta_{1,1,x}\zeta_{0,0,x}}
      {\zeta_{1,0,x}\zeta_{0,1,x}}
 =\frac{w_{1,x}(1-w_{0,x})}{w_{0,x}(1-w_{1,x})}.                       \tag{41}
\]
After cancellation of the positive factor \(\lambda_x^2p_x(1-p_x)\), Assumptions \(\mathrm{ass{:}treatment\text{-}odds\text{-}upper}\) and \(\mathrm{ass{:}treatment\text{-}odds\text{-}lower}\) are exactly \((w_{0,x},w_{1,x})\in D_{A,x}\).

On the strictly positive chart, put \(\theta_{a,x}=\operatorname{OR}(\mu_{a,1,x},\mu_{a,0,x})\).  Cancelling the positive \((U,X=x)\)-margins in the two outcome cross products, Assumptions \(\mathrm{ass{:}outcome\text{-}odds\text{-}upper}\) and \(\mathrm{ass{:}outcome\text{-}odds\text{-}lower}\) give
\[
 \Gamma_a^{-1}\leq\theta_{a,x}\leq\Gamma_a.                          \tag{42}
\]
Lemma \(\mathrm{lem{:}binary\text{-}mixture\text{-}inversion}\) applied to (39) uniquely gives
\[
 \mu_{a,0,x}=\kappa(r_{a,x},w_{a,x},\theta_{a,x}),
 \qquad
 \mu_{a,1,x}
 =\frac{\theta_{a,x}\mu_{a,0,x}}
        {1+(\theta_{a,x}-1)\mu_{a,0,x}},                               \tag{43}
\]
with both risks strictly between zero and one.  Subtracting (39) from the marginal potential-outcome mean and using (40) gives the algebraic identity
\[
\begin{aligned}
 \widetilde E\{Y(a)\mid X=x\}
 &=(1-\pi_x)\mu_{a,0,x}+\pi_x\mu_{a,1,x}\\
 &=r_{a,x}+(\pi_x-w_{a,x})(\mu_{a,1,x}-\mu_{a,0,x})\\
 &=q_{a,x}(w_{0,x},w_{1,x};\theta_{a,x}).
\end{aligned}                                                           \tag{44}
\]
Thus the observed-law functional is linked to the causal mean by the assumptions; it has not been defined to equal it.

Lemma \(\mathrm{lem{:}outcome\text{-}budget\text{-}extrema}\) optimizes (44) over (42), separately in the two outcome arms but at the same posterior pair.  Hence
\[
 H_x^-(w_{0,x},w_{1,x})
 \leq \widetilde E\{Y(1)-Y(0)\mid X=x\}
 \leq H_x^+(w_{0,x},w_{1,x}),                                         \tag{45}
\]
and minimization or maximization over the common feasible pair gives
\[
 \ell_x\leq\widetilde E\{Y(1)-Y(0)\mid X=x\}\leq\upsilon_x.          \tag{46}
\]
After multiplication by \(\lambda_x>0\) and summation over the finite support,
\[
 L\leq\tau(\widetilde P)\leq R,
 \qquad \widetilde P\in\mathcal M^{\circ}.                            \tag{47}
\]
If \(\widetilde P\in\mathcal M\), Definition \(\mathrm{def{:}common\text{-}u\text{-}class}\) gives \(\widetilde P_k\in\mathcal M^{\circ}\) converging to \(\widetilde P\) in the finite simplex.  Since \(\tau\) is linear and continuous, (47) passes to the limit.  Therefore
\[
 \Theta_I(P_{\mathrm{obs}})\subseteq[L,R].                             \tag{48}
\]

We construct compatible laws for the reverse inclusion.  In every stratum choose \((w_{0,x},w_{1,x})\in D_{A,x}\cap(0,1)^2\) and \(\theta_{a,x}\in[\Gamma_a^{-1},\Gamma_a]\), and set
\[
\begin{array}{ll}
 \widetilde P(A=1,U=1\mid X=x)=p_xw_{1,x},
 &\widetilde P(A=1,U=0\mid X=x)=p_x(1-w_{1,x}),\\
 \widetilde P(A=0,U=1\mid X=x)=(1-p_x)w_{0,x},
 &\widetilde P(A=0,U=0\mid X=x)=(1-p_x)(1-w_{0,x}).
\end{array}                                                             \tag{49}
\]
Observed overlap and interiority of the selected pair make all four probabilities positive.  Define \(\mu_{a,u,x}\) by (43).  Lemma \(\mathrm{lem{:}binary\text{-}mixture\text{-}inversion}\) puts them in \((0,1)\) and makes (39) hold.  Conditional on \((X=x,U=u)\), draw \(Y(0)\) and \(Y(1)\) independently as Bernoulli variables with success probabilities \(\mu_{0,u,x}\) and \(\mu_{1,u,x}\); draw \(A\) according to (49), independently of the potential outcomes given \((X,U)\); and set
\[
 Y=AY(1)+(1-A)Y(0).                                                     \tag{50}
\]
This construction proves \((Y(0),Y(1))\perp A\mid(X,U)\).  Equations (39), (49), and (50) reproduce the observed conditional law of \((A,Y)\).  Equation (41), membership in \(D_{A,x}\), and (42) verify all four odds restrictions.  With \(\widetilde P(X=x)=\lambda_x\), every full-data cell is positive, so the constructed law belongs to \(\mathcal M^{\circ}\), and (44) gives its causal means.

It remains to show that this interior construction attains every value in \([L,R]\), including its endpoints.  Define
\[
 \mathcal D_x^{\circ}
 =\{D_{A,x}\cap(0,1)^2\}
  \times[\Gamma_0^{-1},\Gamma_0]
  \times[\Gamma_1^{-1},\Gamma_1].                                    \tag{51}
\]
The coordinatewise logit map sends the posterior-pair factor bijectively to the convex strip
\[
 \{(s_0,s_1)\in\mathbb R^2:
      |s_1-s_0|\leq\log\Gamma_A\},                                    \tag{52}
\]
so that factor, and therefore \(\mathcal D_x^{\circ}\), is connected.  Lemma \(\mathrm{lem{:}binary\text{-}mixture\text{-}inversion}\) and Proposition \(\mathrm{prop{:}closed\text{-}chart\text{-}extension}\) make the stratum contrast continuous on (51).

Suppose first that \(\Gamma_A>1\) and \((\Gamma_0,\Gamma_1)\neq(1,1)\).  Theorem \(\mathrm{thm{:}one\text{-}arc\text{-}endpoint\text{-}reduction}\) gives canonical-arc optimizers over \(0\leq t\leq1\).  At \(t=0\) or \(t=1\), the pair is constant and both endpoint objectives equal \(r_{1,x}-r_{0,x}\).  For any \(0<t<1\), \(W_{\Gamma_A}(t)>t\).  The formulas and strict tilt-gap positivity in that theorem then give
\[
 H_x^+(W_{\Gamma_A}(t),t)>r_{1,x}-r_{0,x},
 \qquad
 H_x^-(W_{\Gamma_A}(t),t)<r_{1,x}-r_{0,x},                             \tag{53}
\]
because at least one of \(\Gamma_0,\Gamma_1\) exceeds one and \(p_x,1-p_x>0\).  Thus neither optimizer is an endpoint of the arc, and compactness supplies optimizing \(t\)'s with \(0<t<1\).  Their posterior pairs lie in (51), and Lemma \(\mathrm{lem{:}outcome\text{-}budget\text{-}extrema}\) supplies the endpoint tilts that attain \(\ell_x\) and \(\upsilon_x\).  If instead \(\Gamma_A=1\) or \(\Gamma_0=\Gamma_1=1\), the one-arc theorem gives \(\ell_x=\upsilon_x=r_{1,x}-r_{0,x}\); this value is attained in (51) by \(w_{0,x}=w_{1,x}=1/2\) and \(\theta_{0,x}=\theta_{1,x}=1\).

Therefore the continuous image of connected \(\mathcal D_x^{\circ}\) lies in \([\ell_x,\upsilon_x]\) by (45) and contains both endpoints, so it equals that interval.  The finite product \(\prod_x\mathcal D_x^{\circ}\) is connected, and the continuous weighted-sum map has image
\[
 \left[\sum_x\lambda_x\ell_x,
       \sum_x\lambda_x\upsilon_x\right]=[L,R].                       \tag{54}
\]
Applying (49)--(50) to every point in this product proves
\[
 [L,R]
 \subseteq\{\tau(\widetilde P):\widetilde P\in\mathcal M^{\circ}\}
 \subseteq\{\tau(\widetilde P):\widetilde P\in\mathcal M\}.
                                                                            \tag{55}
\]
Together with (48), this establishes
\[
 \{\tau(\widetilde P):\widetilde P\in\mathcal M^{\circ}\}
 =\{\tau(\widetilde P):\widetilde P\in\mathcal M\}
 =[L,R].                                                                 \tag{56}
\]
The endpoint laws constructed above use independent conditional Bernoulli potential outcomes.  More generally, any strictly positive coupling of the same two Bernoulli marginals given \((X,U)\) preserves every displayed calculation and conditional ignorability.  A coupling with zero cells belongs to the closed class after mixing it with an arbitrarily small amount of the independent coupling, which preserves the marginals and converges in the finite simplex.  This final closure observation concerns alternative couplings only; equation (56) shows that no endpoint requires a boundary posterior pair.  Rosenbaum--Rubin is used only as a historical fixed-tuple comparator, and no historical class equality, inclusion, or numerical dominance is used.
